\chapter{WP1: Modello}
\label{ch:wp1}

Nei sistemi democratici moderni, la conduzione di elezioni che siano al contempo sicure, trasparenti e verificabili rappresenta un requisito fondamentale. Il voto elettronico si propone come una soluzione capace di incrementare l'accessibilità alle urne, abbattere i costi logistici e accelerare sensibilmente le operazioni di scrutinio. Tuttavia, la transizione verso il digitale introduce criticità significative in termini di sicurezza, affidabilità e fiducia da parte degli elettori.
Le sfide principali riguardano la possibile violazione della segretezza del voto, il rischio di manipolazione o alterazione dei dati memorizzati e l'impossibilità, nei sistemi tradizionali, per il singolo elettore di verificare l'effettivo conteggio della propria preferenza. A differenza del voto fisico, in un sistema digitale la fiducia è riposta in componenti software e hardware che possono però essere soggetti a malfunzionamenti o attacchi mirati.

In questo capitolo viene definito il modello del sistema di voto elettronico proposto. Il Work Package 1 (WP1) ha lo scopo di delineare la funzionalità che si intende realizzare, individuando gli attori coinvolti e i loro obiettivi. 
Verrà inoltre presentato un \textit{threat model} dettagliato, volto a identificare i potenziali avversari che potrebbero agire contro il sistema, le risorse a cui hanno accesso e i tipi di attacchi che potrebbero tentare di eseguire. Questo processo è fondamentale per comprendere le minacce a cui il protocollo può essere esposto e per guidare la progettazione delle contromisure necessarie a garantire la sicurezza del sistema. 
Infine, verranno formalizzate le proprietà di resilienza che il protocollo deve preservare in presenza di attacchi, stabilendo i criteri di valutazione per la progettazione e l'analisi dei successivi Work Package. 

La definizione accurata del modello costituisce la base necessaria per lo sviluppo di un sistema di voto elettronico che sia non solo funzionale, ma anche sicuro e affidabile, in grado di resistere a una vasta gamma di minacce e di garantire la fiducia degli elettori e degli osservatori.

\clearpage


\section{Tipologia di elezione}
\label{sec:1.1_tipologia_elezione}

Per questo sistema si è stabilito di operare su un modello di \textbf{elezione a lista chiusa} caratterizzato da un numero $N$ generico di opzioni di voto. In tale configurazione, l'elettore esprime la propria preferenza selezionando una singola opzione all'interno di un insieme prefissato, senza la possibilità di indicare preferenze nominali libere o di ordinare i candidati secondo un sistema di voto preferenziale. 
Il sistema può essere ricondotto anche al caso particolare di referendum binario, nel quale l'insieme delle opzioni di voto è costituito unicamente da due alternative possibili (ad esempio ``Sì'' e ``No''). Tale generalizzazione permette di mantenere il modello sufficientemente astratto da poter essere applicato a differenti scenari elettorali.

La scelta di questo specifico modello, in conformità con i requisiti di progetto, è motivata da specifiche implicazioni sulla complessità strutturale del protocollo e sui compromessi di sicurezza adottabili. Dal punto di vista crittografico e computazionale, limitare il dominio delle scelte a un insieme discreto, finito e noto a priori riduce drasticamente la complessità delle operazioni necessarie a garantire la validità del voto. La gestione di una preferenza nominale aperta richiederebbe il trattamento di input a lunghezza variabile, esponendo il sistema a problematiche complesse di validazione e aumentando sensibilmente il carico computazionale.  

Al contrario, l'utilizzo di una lista chiusa permette di modellare i voti attraverso un formato logico predefinito e rigorosamente omogeneo. Questo approccio astrae e semplifica il soddisfacimento dei requisiti di confidenzialità e integrità: garantendo che ogni singola scheda appartenga a un dominio ristretto, si facilita enormemente la verifica della sua conformità strutturale.

Il compromesso fondamentale adottato in questa scelta progettuale risiede quindi nel privilegiare l'efficienza algoritmica e l'integrità a discapito della massima generalità del tipo di voto. Si ritiene che, per gli scopi di questo protocollo, garantire in modo incondizionato la correttezza del conteggio e la robustezza strutturale delle schede sia prioritario e strategicamente più solido rispetto all'implementazione di sistemi di preferenza complessi.

\clearpage


\section{Funzionalità del sistema}
\label{sec:1.2_funzionalità_sistema}

La funzionalità primaria del protocollo consiste nella gestione di un'elezione elettronica sicura all'interno di un ambiente remoto e non fisicamente controllato, garantendo che l'intero processo sia trasparente, verificabile e protetto da manipolazioni sistemiche. Al fine di delineare in modo dettagliato e non ambiguo il perimetro logico del sistema, il funzionamento del protocollo viene modellato attraverso un flusso sequenziale suddiviso in tre macro-fasi funzionali e cronologiche, propedeutiche l'una all'altra:

\begin{enumerate}
    \item \textbf{Fase di Registrazione e Autenticazione:} In questa fase iniziale, il sistema ha il compito di verificare l'eleggibilità dell'utente che richiede di accedere alla votazione. Il perimetro funzionale impone l'accertamento che l'utente sia legittimamente incluso nelle liste degli aventi diritto e che non abbia ancora speso il proprio diritto di voto. L'output logico di questa fase è l'emissione di un'autorizzazione formale che abiliti l'utente alle fasi successive, garantendo tassativamente che tale autorizzazione non contenga elementi idonei a svelare l'identità reale del richiedente.
    \item \textbf{Fase di Votazione ed Espressione del Suffragio:} Una volta ottenuta l'abilitazione, l'elettore predispone la propria scheda selezionando una singola preferenza all'interno del dominio ristretto della lista chiusa. Il sistema garantisce la sottomissione sicura della scheda in modalità confidenziale, inserendola all'interno di un registro delle transazioni elettorali. Come output di questa fase, il sistema genera e restituisce una ricevuta logica (es. un codice di tracciamento della scheda cifrata) dell'avvenuta corretta sottomissione, che fungerà da elemento di controllo per la successiva fase di verifica.
    \item \textbf{Fase di Scrutinio e Pubblicazione:} Al raggiungimento del termine temporale stabilito, il sistema impedisce in modo definitivo la ricezione di qualunque ulteriore scheda, dichiarando chiusa la raccolta. Viene quindi avviata la procedura di calcolo aggregato delle preferenze registrate nel sistema. La funzionalità richiesta impone che l'aggregazione dei risultati avvenga in modo matematicamente separato rispetto alle singole schede, terminando con la pubblicazione del conteggio finale dei voti e delle relative prove di correttezza necessarie alla validazione indipendente da parte di osservatori esterni.
\end{enumerate}

\subsection*{Vincoli operativi e Regole del voto}
Il modello definisce due vincoli funzionali necessari per preservare la correttezza formale del processo:
\begin{itemize}
    \item \textbf{Vincolo di Unicità (One-Man, One-Vote):} Ciascun elettore ha diritto a esprimere un solo voto valido nel conteggio finale. Il sistema deve riconoscere e bloccare automaticamente qualsiasi tentativo di inserimento multiplo o ridondante associato al medesimo diritto di voto.
    \item \textbf{Non-correlabilità dei Flussi Informativi:} Il perimetro logico del sistema impone una barriera strutturale invalicabile tra i dati crittografici utilizzati per validare l'identità e l'eleggibilità dell'utente e i dati crittografici che contengono la preferenza espressa. Questo disaccoppiamento dei flussi di informazione è un requisito funzionale primario per garantire matematicamente che l'omogeneità delle schede impedisca la ricostruzione a ritroso del legame tra votante e voto.
\end{itemize}

\subsection*{Analisi sulla Revocabilità e Sostituzione del Voto}
In conformità con le direttive di progetto, è stata attentamente valutata l'opportunità di consentire all'elettore di modificare, sovrascrivere o annullare il proprio voto prima della chiusura formale delle urne. Questa specifica proprietà introduce un profondo conflitto strutturale tra la resilienza alla coercizione, lo pseudo-anonimato e l'efficienza complessiva del sistema.

Se da un lato la revocabilità offre una difesa teorica contro il voto di scambio e la coercizione in ambienti remoti non controllati (permettendo a un elettore coartato di votare sotto minaccia e successivamente esprimere il vero voto in isolamento), dall'altro lato la sua introduzione comporta criticità sistemiche che in questo protocollo vengono ritenute inaccettabili per le seguenti motivazioni:
\begin{enumerate}
    \item \textbf{Minaccia allo Pseudo-anonimato:} Per consentire la sostituzione o l'aggiornamento di una scheda precedentemente inviata, l'infrastruttura di memorizzazione dovrebbe necessariamente disporre di un meccanismo per associare le diverse transazioni espresse dallo stesso utente (ad esempio tramite un identificativo persistente), oppure richiedere all'elettore di dimostrare la paternità del voto precedente. Questo legame logico aumenta la superficie d'attacco, innalzando il rischio di violazione definitiva dell'anonimato.
    \item \textbf{Violazione del Vincolo di Efficienza:} Gestire la purificazione di voti multipli o la sovrascrittura dinamica delle schede richiede un'architettura logica complessa e protocolli di interazione continui. Ciò incrementerebbe in modo esponenziale la latenza di rete e il costo computazionale, violando il vincolo di efficienza che costituisce un limite progettuale stringente fin da questo primo Work Package.
    \item \textbf{Complessità della Verificabilità:} Permettere voti modificabili complicherebbe le procedure di verifica individuale e universale. Gli osservatori non dovrebbero semplicemente verificare l'inclusione di una scheda in un registro immutabile, ma dovrebbero validare la corretta cronologia e l'effettivo scarto delle schede obsolete, richiedendo strumenti logici avanzati che esulano dagli obiettivi del sistema.
\end{enumerate}

Sulla base di questa onesta analisi critica dei limiti strutturali, si stabilisce come scelta progettuale definitiva che \textbf{il voto espresso sia rigidamente non revocabile e non modificabile (modello "one-shot")}. Una volta che la scheda viene trasmessa e registrata, essa è considerata definitiva e immutabile. Il rischio residuo legato alla coercizione estemporanea viene accettato e mitigato parzialmente attraverso la rigidità del formato a lista chiusa analizzato nella sezione precedente, privilegiano in modo assoluto la linearità del protocollo, la robustezza dello pseudo-anonimato e la massima efficienza computazionale globale del sistema.


\section{Attori Onesti del Sistema}
\label{sec:1.3_attori_onesti}

L'architettura logica del protocollo si fonda sull'interazione di quattro entità funzionali distinte, ciascuna caratterizzata da uno specifico stato informativo, compiti operativi definiti e precisi obiettivi di sicurezza da conseguire. La divisione delle funzioni presentata di seguito è finalizzata a distribuire le responsabilità operative e a prevenire la concentrazione di poteri computazionali o informativi in un singolo punto del sistema. Ciascun attore opera in modo indipendente, garantendo che il corretto svolgimento delle elezioni sia il risultato della cooperazione di componenti distinte.

L'\textbf{Elettore} è l'utente che ha il diritto di partecipare alla consultazione. All'atto pratico, il suo compito è selezionare una sola opzione tra le $N$ liste chiuse disponibili, producendo un voto valido che impedisca scelte multiple o irregolari. L'Elettore ha due obiettivi principali: da una parte, vuole la garanzia che il suo voto resti completamente segreto, senza che nessuno possa mai collegare la sua identità alla preferenza espressa. Dall'altra, ha bisogno di esercitare la verificabilità individuale. Deve cioè poter controllare in autonomia che la sua scheda sia finita intatta nell'urna elettronica e sia stata conteggiata nel risultato finale. Questa verifica, però, deve essere progettata in modo tale da garantire la proprietà di Receipt-Freeness (assenza di ricevuta): il sistema non deve fornirgli alcuna prova crittografica che possa essere esibita a terzi per dimostrare inequivocabilmente il contenuto della sua preferenza.

Il \textbf{Sistema di Autenticazione (SA)} agisce come un Identity Provider (IdP) e ha il compito principale di gestire le identità e verificare che i partecipanti abbiano i requisiti necessari per votare. Questa componente si occupa di validare gli utenti legittimi prima dell'inizio delle operazioni e di vigilare sul rispetto del principio di unicità del voto, assicurando che ogni avente diritto possa esprimere una sola preferenza all'interno del sistema. Per preservare lo pseudo-anonimato, l'attività del Sistema di Autenticazione deve essere separata in modo netto dalla successiva fase di raccolta delle schede, prevedendo che le due funzioni siano gestite da entità del tutto indipendenti. Tale disaccoppiamento garantisce che la procedura di verifica dell'identità avvenga all'oscuro della reale preferenza espressa dall'utente, impedendo a chi gestisce le credenziali di tracciare il momento esatto della sottomissione o di collegare l'anagrafica dell'elettore al voto inserito nell'urna.

L'\textbf{Autorità Elettorale} è l'entità il cui compito è quello di ricevere, ordinare e conservare le schede protette all'interno dell'urna elettronica. Quando le votazioni si chiudono, l'Autorità Elettorale deve eseguire lo scrutinio e calcolare i voti totali ottenuti da ciascuna delle $N$ liste. Questa operazione deve avvenire attraverso procedure che permettono di aggregare i risultati mantenendo i singoli voti del tutto coperti. In questo modo, il conteggio finale viene calcolato senza che nessuna entità possa collegare l'identità dell'elettore alla preferenza specifica memorizzata nell'urna.

Il \textbf{Verificatore Pubblico} è un attore indipendente e passivo, che può essere un osservatore esterno, un comitato di controllo o l'elettore stesso. Il suo unico scopo è verificare che le operazioni di voto e di scrutinio si svolgano in modo corretto. Questa entità non partecipa alla creazione dei voti e non si occupa di raccoglierli, ma interviene esaminando le prove e i dati pubblici messi a disposizione dall'Autorità Elettorale. La presenza del Verificatore Pubblico è indispensabile per realizzare la verificabilità universale, dando a chiunque la possibilità di accertarsi, in totale autonomia, che il risultato finale pubblicato sia l'esatto calcolo delle schede memorizzate nell'urna. L'inserimento di questa figura nel modello è fondamentale per garantire la trasparenza del processo elettorale e per rafforzare la fiducia degli elettori nei confronti del sistema di voto elettronico adottato.

\section{Assunzioni di Fiducia}
\label{sec:1.4_assunzioni_fiducia}

La sicurezza di un sistema di voto elettronico non è mai assoluta, ma è intrinsecamente legata a un insieme di ipotesi sul comportamento delle entità partecipanti e sull'ambiente operativo. Di seguito vengono analizzate le assunzioni di fiducia fondamentali per il modello proposto, spiegando perché le abbiamo scelte e il rischio derivante da una loro potenziale violazione. In generale, il sistema garantisce le proprietà di segretezza, integrità e verificabilità unicamente sotto l'ipotesi che le seguenti condizioni siano verificate.

\subsection*{1. Non collusione tra Sistema di Autenticazione e Autorità Elettorale}
\textbf{Dichiarazione:} Si assume che il Sistema di Autenticazione (SA) e l'Autorità Elettorale (AE) operino in modo del tutto indipendente, senza scambiarsi informazioni di nascosto. \\
\textbf{Motivazione e compromessi:} Il sistema garantisce lo pseudo-anonimato a condizione che tale separazione venga rispettata. Il SA sa chi vota, mentre l'AE sa cosa viene votato. Tenerli separati è un compromesso pratico per avere un sistema leggero ed efficiente, fidandoci del fatto che i due gestori non si accordino per imbrogliare. \\
\textbf{Rischio Residuo:} Se questa assunzione crolla e le due entità comunicano, la garanzia di pseudo-anonimato viene totalmente compromessa. Gli amministratori potrebbero facilmente incrociare i dati di accesso con l'arrivo delle schede, ricostruendo l'esatto legame tra l'identità anagrafica dell'elettore e la preferenza depositata.

\subsection*{2. Integrità del terminale dell'Elettore}
\textbf{Dichiarazione:} Si assume che il dispositivo usato dall'Elettore non sia compromesso da malware, keylogger o spyware. \\
\textbf{Motivazione e compromessi:} Il sistema garantisce la segretezza e la correttezza del singolo voto a condizione che l'ambiente di esecuzione locale sia affidabile. Votando da remoto, l'assenza di un ambiente fisico controllato obbliga a delegare la sicurezza del dato in chiaro al dispositivo dell'elettore. Purtroppo, il sistema non può difendersi da minacce che agiscono direttamente sullo schermo o sulla tastiera dell'utente. \\
\textbf{Rischio Residuo:} In caso di dispositivo infetto, un attaccante software potrebbe leggere la scelta elettorale prima ancora che venga protetta, violando la segretezza. Nel peggiore dei casi, il malware potrebbe manipolare l'input, inviando all'urna una preferenza diversa dalla volontà reale dell'elettore.

\subsection*{3. Immutabilità dell'urna elettronica}
\textbf{Dichiarazione:} Si assume che l'Autorità Elettorale conservi i dati onestamente, funzionando come un'urna in cui si possono solo aggiungere schede, senza poterne mai cancellare o alterare nessuna. \\
\textbf{Motivazione e compromessi:} Il sistema assicura l'integrità del conteggio a condizione che l'AE non elimini i record. Affidare i dati a un server centrale semplifica decisamente la raccolta e lo scrutinio dei voti, ma ci costringe a fidarci del fatto che l'amministratore non cancelli di proposito le schede di una lista avversaria. \\
\textbf{Rischio Residuo:} Se l'AE procede di nascosto alla cancellazione di alcuni voti, il risultato finale sarà falsato. Questo rischio è in parte mitigato perché l'elettore non troverebbe più la sua scheda controllando l'urna pubblica, ma il sistema in sé non può impedire fisicamente all'amministratore di effettuare la cancellazione.

\subsection*{4. Spazio fisico sicuro e assenza di coercizione}
\textbf{Dichiarazione:} Si assume che l'Elettore possa votare in un momento e in un luogo privato, senza subire minacce o ricatti fisici. \\
\textbf{Motivazione e compromessi:} In questo progetto abbiamo scelto di non permettere la modifica o la revoca del voto. Il sistema garantisce un conteggio certo e inequivocabile a condizione che il voto inviato sia definitivo. È una scelta che rende il protocollo molto più semplice da gestire ed efficiente, ma rinuncia a uno strumento utile contro la compravendita dei voti. \\
\textbf{Rischio Residuo:} Non potendo cambiare la scelta in un secondo momento, se qualcuno costringesse con la forza un elettore a votare davanti a lui, il ricatto andrebbe a buon fine. Il rischio residuo è la totale perdita di libertà del votante, un limite molto difficile da superare nei sistemi di voto a distanza.

\subsection*{5. Sicurezza e affidabilità delle primitive crittografiche}
\textbf{Dichiarazione:} Si assume che gli algoritmi crittografici scelti siano sicuri e che nessun attaccante abbia a disposizione una potenza di calcolo sufficiente per comprometterli in tempi brevi. \\
\textbf{Motivazione e compromessi:} Il sistema garantisce la segretezza e l'inviolabilità delle schede a condizione che la matematica alla base delle funzioni scelte sia solida. Questa è un'ipotesi inevitabile per qualsiasi software moderno: non si può progettare un protocollo digitale senza fidarsi degli strumenti che proteggono i dati. Il compromesso è che accettiamo un rischio teorico legato al futuro progresso tecnologico pur di poter implementare il sistema oggi. \\
\textbf{Rischio Residuo:} Se un attaccante trovasse un modo per aggirare l'algoritmo crittografico, tutte le difese del sistema crollerebbero. Chi attacca potrebbe calcolare le chiavi segrete del server, riuscendo a leggere i voti individuali degli utenti e persino a modificare i risultati finali dello scrutinio senza lasciare alcuna traccia.


\section{Modello di Minaccia (Threat Model)}
\label{sec:1.5_threat_model}

Definire il Modello di Minaccia serve a inquadrare in modo realistico le entità ostili a cui il sistema deve saper resistere. Nel rispetto delle assunzioni di fiducia appena descritte abbiamo identificato diversi profili di attacco distinti. È importante sottolineare che l'inclusione di attori interni come il Sistema di Autenticazione o l'Autorità Elettorale non contraddice le assunzioni precedenti, ma serve a valutare la robustezza del protocollo nel caso in cui un singolo componente devii dal comportamento previsto. Ognuno di questi avversari dispone di mezzi differenti, persegue scopi mirati e colpisce una specifica area del protocollo.

\subsection*{1. Avversario di Rete (Man-in-the-Middle)}
\textbf{Tipologia:} Attivo / Passivo. \\
\textbf{Descrizione:} Questo avversario opera direttamente sul canale di trasmissione dei dati. In modalità passiva, tenta di violare la segretezza spiando i pacchetti in transito per scoprire l'associazione tra elettori e liste scelte. In modalità attiva, il suo obiettivo è alterare la comunicazione in transito, cercando di modificare il voto contenuto all'interno di un pacchetto prima che questo raggiunga il server dell'urna elettronica. \\
\textbf{Risorse:} Dispone del controllo totale dell'infrastruttura di rete, comprese le connessioni internet pubbliche, le reti locali e i nodi dei provider di servizi. Questa posizione strategica gli permette di leggere, bloccare, ritrasmettere o tentare di manipolare qualsiasi dato scambiato tra i vari attori del sistema, pur non possedendo l'accesso ad alcuna chiave crittografica segreta.

\subsection*{2. Attaccante DoS (Denial of Service)}
\textbf{Tipologia:} Attivo. \\
\textbf{Descrizione:} L'obiettivo di questo profilo è bloccare la piattaforma elettorale per impedirne il corretto utilizzo. Non mira a rubare informazioni o a modificare i dati interni, ma vuole compromettere la disponibilità del servizio, facendo in modo che gli utenti non riescano a raggiungere i server o a inviare la propria preferenza prima della scadenza del tempo. \\
\textbf{Risorse:} Ha a disposizione una notevole larghezza di banda e strumenti capaci di generare enormi flussi di traffico di rete. Utilizza queste risorse per saturare le connessioni verso i server dedicati all'autenticazione e alla raccolta delle schede, fino a provocare il blocco delle comunicazioni.

\subsection*{3. Falsario}
\textbf{Tipologia:} Attivo. \\
\textbf{Descrizione:} Questo avversario si concentra sulla manipolazione dei meccanismi di validazione. Il suo obiettivo è ingannare il sistema sottomettendo schede elettorali volutamente malformate oppure falsificando le firme digitali per forzare l'accettazione di voti illegittimi all'interno dell'urna, eludendo di fatto i controlli di integrità. \\
\textbf{Risorse:} Dispone di notevoli capacità di calcolo e di una conoscenza profonda degli algoritmi del sistema. Interagisce con il protocollo dall'esterno con lo scopo di aggirare i controlli di sicurezza del sistema, ma non ha alcun accesso amministrativo ai server.

\subsection*{4. Elettore Disonesto (o Impostore)}
\textbf{Tipologia:} Attivo. \\
\textbf{Descrizione:} È un utente che tenta di violare i vincoli e le procedure del protocollo elettorale. I suoi obiettivi spaziano dal tentare di votare più volte, all'usare le credenziali rubate a un altro studente per votare al suo posto. Tenta inoltre di estrarre dal sistema una "ricevuta" matematica da poter vendere a terzi per dimostrare inequivocabilmente come ha votato (voto di scambio). \\
\textbf{Risorse:} Possiede credenziali di accesso valide e conosce il funzionamento dell'interfaccia di voto ma, come il falsario, non ha alcun potere gestionale sull'infrastruttura centrale.

\subsection*{5. Il Sistema di Autenticazione Disonesto}
\textbf{Tipologia:} Attivo / Passivo. \\
\textbf{Descrizione:} Questo avversario sfrutta il proprio ruolo di controllo degli accessi per minare la riservatezza delle elezioni e la correttezza della partecipazione. Quando agisce in modo passivo, cerca di ricostruire il legame tra l'identità dell'utente e la preferenza espressa, provando a incrociare i dati di chi si collega con i voti che appaiono pubblicamente. Quando agisce in modo attivo, tenta di alterare l'andamento del voto escludendo ingiustamente utenti che avrebbero il diritto di partecipare o, al contrario, permettendo il voto a soggetti non autorizzati. \\
\textbf{Risorse:} Opera dall'interno del sistema con il controllo completo sulla gestione delle identità e sui permessi di ingresso alla piattaforma. Ha la capacità di tracciare le richieste di accesso e di emettere le autorizzazioni necessarie per votare, ma non ha alcun potere di gestione sull'urna elettronica né la possibilità di vedere o manipolare le chiavi destinate al conteggio dei voti.

\subsection*{6. Autorità Elettorale Disonesta}
\textbf{Tipologia:} Attivo / Passivo. \\
\textbf{Descrizione:} Questo avversario sfrutta il suo ruolo centrale nella raccolta delle schede per falsare l'esito della consultazione. Quando agisce in modo passivo, cerca di scoprire le preferenze in anticipo per monitorare segretamente l'andamento delle votazioni. Quando opera in modo attivo, tenta di alterare il conteggio finale, scartando di proposito voti validi oppure inserendo preferenze arbitrarie all'interno dell'urna. \\
\textbf{Risorse:} Rappresenta la minaccia interna con il controllo più ampio sui dati raccolti. Gestisce direttamente l'infrastruttura che ospita l'urna digitale e la bacheca pubblica accessibile a tutti. Ha in custodia i voti protetti e si occupa materialmente del calcolo dei risultati, ma non ha a disposizione le informazioni necessarie per collegare una specifica scheda alla persona reale che l'ha inviata.


\clearpage


\section{Proprietà di Sicurezza da Preservare}
\label{sec:1.6_proprieta_sicurezza}
% Declinazione delle garanzie di sicurezza in "micro-proprietà" atomiche, per facilitare la successiva argomentazione analitica.  
% Suddivisione in macro-categorie come da traccia:
% Autenticità e Unicità .  
% Confidenzialità e Segretezza.  
% Integrità e Invarianza.  
% Verificabilità (Individuale e Universale)

A fronte del Modello di Minaccia delineato nella sezione precedente, il protocollo crittografico è progettato per garantire un set di proprietà di sicurezza fondamentali. Tali proprietà definiscono le capacità di resilienza del sistema in presenza di attacchi attivi o passivi, assicurando che gli obiettivi primari dell'elezione vengano preservati o, nell'eventualità di una compromissione, che l'attacco venga matematicamente rilevato. 

Di seguito vengono formalizzate le proprietà che il sistema si prefigge di garantire, esplicitando per ciascuna di esse la minaccia che mira a neutralizzare.

\subsection*{1. Autenticità}
\textbf{Definizione:} Solo ed esclusivamente gli elettori legittimi, preventivamente identificati e autenticati tramite credenziali crittografiche verificabili, hanno il diritto e la capacità tecnica di partecipare alla consultazione elettorale.
\\
\textbf{Resilienza agli attacchi:} Questa proprietà costituisce la difesa primaria contro il \textit{Falsario} e l'\textit{Elettore Disonesto (Impostore)}. Il sistema garantisce che nessun attaccante esterno possa iniettare schede valide nell'urna senza possedere un'abilitazione regolarmente emessa. Inoltre, previene lo \textit{spoofing} dell'identità, impedendo a un utente di votare a nome di un altro elettore senza esserne in possesso delle credenziali crittografiche.

\subsection*{2. Unicità (One-Man, One-Vote)}
\textbf{Definizione:} Ciascun elettore legittimo può far confluire al massimo un solo voto valido nel conteggio finale dello scrutinio.
\\
\textbf{Resilienza agli attacchi:} Questa proprietà contrasta direttamente l'\textit{Elettore Disonesto} che tenta di sbloccare o aggirare i limiti dell'interfaccia per sottomissioni multiple (doppio voto). La garanzia di unicità si fonda sul corretto funzionamento del Sistema di Autenticazione, il quale deve tracciare l'avvenuta partecipazione dell'utente senza tuttavia apporre firme identificabili sulla scheda cifrata, mantenendo così il disaccoppiamento logico.

\subsection*{3. Segretezza e Pseudo-anonimato}
\textbf{Definizione:} Il contenuto del voto espresso dal singolo elettore non deve in alcun modo essere rivelabile o associabile alla sua identità reale. Tale garanzia deve resistere anche a tentativi di de-anonimizzazione basati sull'incrocio delle informazioni di traffico o dei metadati.
\\
\textbf{Resilienza agli attacchi:} La segretezza protegge l'elettore dall'\textit{Avversario di Rete (passivo)}, che tenta di intercettare le schede in chiaro, e dalle entità interne (\textit{Sistema di Autenticazione Disonesto} e \textit{Autorità Elettorale Disonesta} in modalità passiva). Sotto l'assunzione di non-collusione (Sezione \ref{sec:1.4_assunzioni_fiducia}), l'architettura a ruoli separati garantisce che né chi conosce l'identità (SA) né chi calcola i voti (AE) disponga dei dati sufficienti per ricostruire il legame logico tra elettore e preferenza. Inoltre, per scoraggiare l'acquisto dei voti da parte di terzi coercitori, il sistema tutela la segretezza garantendo la proprietà di \textit{Receipt-Freeness}: l'utente non possiede alcuna prova matematica in grado di dimostrare inequivocabilmente il contenuto del proprio voto.

\subsection*{4. Integrità}
\textbf{Definizione:} Le schede elettorali, una volta generate sul dispositivo dell'elettore e trasmesse all'infrastruttura di raccolta, non possono subire alterazioni, modifiche o soppressioni durante il transito o dopo la loro registrazione formale nell'urna elettronica.
\\
\textbf{Resilienza agli attacchi:} L'integrità crittografica vanifica i tentativi dell'\textit{Avversario di Rete (attivo)} di manipolare i pacchetti in volo (man-in-the-middle). Inoltre, impone all'\textit{Autorità Elettorale Disonesta} l'utilizzo di una bacheca pubblica in modalità \textit{append-only}; qualsiasi tentativo di cancellare o alterare un record precedentemente registrato invaliderà le firme crittografiche o le strutture dati, rendendo la manomissione immediatamente evidente.

\subsection*{5. Verificabilità Individuale}
\textbf{Definizione:} Ogni elettore ha a disposizione gli strumenti logici per controllare in totale autonomia che la propria scheda sia stata ricevuta intatta, registrata nell'urna pubblica e correttamente predisposta per essere inclusa nel conteggio finale.
\\
\textbf{Resilienza agli attacchi:} È lo strumento principale fornito all'utente onesto per rilevare operazioni fraudolente da parte di un'\textit{Autorità Elettorale Disonesta (attiva)}. Sfruttando la ricevuta di sottomissione (codice di tracciamento), l'elettore verifica che la propria transazione non sia stata vittima di \textit{data dropping} (cancellazione silenziosa del voto) o intercettata da un attacco \textit{Denial of Service} prima di raggiungere il server.

\subsection*{6. Verificabilità Universale}
\textbf{Definizione:} Il risultato finale dell'elezione, unitamente a tutte le schede cifrate raccolte nell'urna, è corredato da prove matematiche che permettono a chiunque (elettori, auditor esterni, comitati di controllo) di validare pubblicamente ed in modo indipendente la correttezza algoritmica dello scrutinio.
\\
\textbf{Resilienza agli attacchi:} Questa proprietà sottrae la fiducia incondizionata al gestore del sistema. Se l'\textit{Autorità Elettorale Disonesta} tentasse di aggregare i voti in modo fraudolento (aggiungendo voti dal nulla o alterando la somma a favore di una specifica lista), le prove crittografiche prodotte a corredo dello scrutinio fallirebbero. La verificabilità universale garantisce che una frode sul conteggio sia matematicamente impossibile da nascondere agli occhi di un Verificatore Pubblico.


\clearpage


\section{Completeness}
\label{sec:1.7_completeness}

La proprietà di completezza descrive la garanzia di funzionamento del protocollo nel cosiddetto "percorso felice" (\textit{happy path}). Essa stabilisce che, qualora tutte le entità partecipanti (Elettore, Sistema di Autenticazione, Autorità Elettorale e Verificatore Pubblico) si comportino in modo strettamente onesto seguendo fedelmente le regole del protocollo, e assumendo valide le ipotesi di fiducia delineate nella Sezione \ref{sec:1.4_assunzioni_fiducia}, il sistema terminerà con successo producendo un risultato matematicamente esatto e pubblicamente verificabile.

In questo scenario ideale, privo di avversari attivi o passivi, la completezza del modello si articola nelle seguenti garanzie logiche sequenziali:

\begin{enumerate}
    \item \textbf{Garanzia di Accesso:} Ogni elettore legittimo che si interfacci correttamente con il Sistema di Autenticazione otterrà con successo l'abilitazione al voto. Il sistema non genererà falsi negativi, permettendo all'utente di procedere regolarmente verso la fase di votazione.
    
    \item \textbf{Garanzia di Accettazione e Ricevuta:} Se l'elettore onesto compila la propria scheda conformemente al formato previsto (selezionando un'unica opzione tra le $N$ liste chiuse), la scheda supererà il controllo di conformità strutturale. L'Autorità Elettorale accetterà la sottomissione, la registrerà nell'urna elettronica pubblica (in cui è possibile solo aggiungere elementi e non cancellarli) e l'elettore riceverà una ricevuta di sottomissione pienamente valida ai fini della verificabilità individuale.
    
    \item \textbf{Correttezza del Computo Aggregato:} Alla chiusura formale delle urne, l'Autorità Elettorale eseguirà l'algoritmo di scrutinio sui dati raccolti. La proprietà di completezza assicura che il totale calcolato per ciascuna lista corrisponderà esattamente alla corretta aggregazione delle reali preferenze inviate dagli elettori onesti, senza alcuna dispersione o alterazione dei dati.
    
    \item \textbf{Convergenza della Verificabilità (con probabilità 1):} Poiché tutte le schede sono state generate e processate rispettando il protocollo, le prove matematiche prodotte insieme ai risultati dello scrutinio risulteranno intrinsecamente corrette. Di conseguenza, l'algoritmo di verifica universale eseguito dal Verificatore Pubblico sui dati dell'urna restituirà sempre esito positivo. Analogamente, la verifica individuale condotta dall'elettore sulla propria ricevuta andrà a buon fine con certezza assoluta (probabilità pari a 1).
\end{enumerate}

La formalizzazione di questa proprietà è un requisito propedeutico per le successive fasi di modellazione: dimostra infatti la validità algoritmica e strutturale del sistema prima ancora di introdurre il modello di minaccia e di analizzare il comportamento del protocollo in presenza di entità malevole.